\Needspace{15\baselineskip}\fancyhead[L]{\small JUNO. ELEMENTE VII.}
\centerline{JUNO. ELEMENTE VII.}

\begin{center}
\rule{0.5\linewidth}{0.4pt}\\[2pt]
Göttingische gelehrte Anzeigen. Stück 14. S. 129, 130. 23. Januar 1808.\\[2pt]
\rule{0.5\linewidth}{0.4pt}
\end{center}

\par\indent Von der \emph{Juno} sind aus dem vergangenen Jahr nur die Beobachtungen bekannt geworden, die Inspektor \textsc{Bessel} in Lilienthal im März, April und Mai mit dem Kreismikrometer gemacht hat. Wahrscheinlich war dieser Planet so lichtschwach, dass er sich nirgends mit feststehenden Instrumenten beobachten ließ. Überdies waren die meisten Astronomen, solange er sichtbar war, so sehr mit der neu entdeckten \emph{Vesta} beschäftigt, dass Juno darüber etwas vernachlässigt worden sein mag. Unser Professor \textsc{Gauss} hat daher nicht gezögert, jene Beobachtungen zur Verbesserung der bisherigen Elemente der Juno zu verwenden, zumal Herr \textsc{Bessel} mit
\SourceLeaf{298}einem ausgezeichneten Teleskop beobachtet hat und im Umgang mit dem Kreismikrometer besonders geschickt ist. Zunächst wurde mithilfe der Beobachtungen vom Mai die Opposition bestimmt. Dabei ergab sich folgendes Ergebnis:

\[
\begin{array}{l}
1807.\ \text{Mai}\ 17.\quad 12^{\text{u}}\,27^{\text{m}}\,41^{\text{s}}\ \text{m.\,Z. in Göttingen}\\[2pt]
\text{wahre Länge}\ .\ .\ .\ .\ .\ 236^\circ\ 4'\ 57''6.\\[2pt]
\text{wahre geozentrische Breite}\quad 16\;56\;0.1\ \text{nördl.}
\end{array}
\]\par\indent Aus der Verbindung dieser Opposition mit denen von 1804 und 1806 und mit anderen Positionen von 1804, 1805 und 1806 wurden dann die folgenden neuen Elemente (VII.) berechnet. Sie geben sämtliche zugrunde gelegten Beobachtungen möglichst genau wieder. [Sie stehen in diesem Abdruck auf Seite 294 als Teil des Briefes an \textsc{Bode} vom 24. Januar 1808.]\par\indent Die Ephemeride für die nächste Erscheinung der Juno nach diesen neuen Elementen ist bereits fertig, und die Astronomen werden sie in der Monatlichen Correspondenz des Freiherrn \textsc{von Zach} finden. Diese Rechnung war umso notwendiger, als die von Professor \textsc{Bode} im Jahrbuch für 1810 nach den VI. Elementen berechnete Ephemeride die Rektaszensionen um die Zeit der nächsten Opposition um mehr als $\frac{3}{4}$ Grad zu groß angibt. Die auch diesmal noch sehr geringe Lichtstärke der Juno wird das Auffinden und Beobachten sehr erschweren.
